\documentclass[preview]{standalone}

\usepackage{amsmath}
\usepackage{amssymb}
\usepackage{stellar}
\usepackage{definitions}

\begin{document}

\id{psicologia-sviluppo-intelligenza}
\genpage

\section{Sviluppo e basi dell'intelligenza}

\begin{snippetdefinition}{disabilita-intellettiva-definition}{Disabilità intellettiva}
    La \emph{disabilità intellettiva} è una condizione caratterizzata da
    significative limitazioni nel funzionamento intellettivo e nel comportamento
    adattivo, espresso in abilità concettuali, sociali e pratico-adattive. Ha
    origine prima dei \(18\) anni.
\end{snippetdefinition}

\begin{snippet}{disabilita-intellettiva-valutazione}
    Si parla di disabilità intellettiva quando il valore del \(QI\) scende al di
    sotto di due deviazioni standard rispetto alla media. Nel DSM-5 sono
    distinti diversi livelli di gravità, ma la valutazione clinica considera
    anche adattamento sociale, emotivo, familiare e scolastico.
\end{snippet}

\begin{snippetdefinition}{plusdotazione-definition}{Plusdotazione}
    La \emph{plusdotazione}, o \emph{giftedness}, è una condizione in cui il
    quoziente intellettivo supera di circa due deviazioni standard la media
    generale e si accompagna al raggiungimento di traguardi di eccellenza nei
    settori di attività dell'individuo.
\end{snippetdefinition}

\begin{snippet}{plusdotazione-ambiente}
    Un \(QI\) elevato non dipende soltanto da basi genetiche ed ereditarie, ma
    anche da componenti educativo-ambientali. Livello di educazione dei genitori,
    prestigio occupazionale familiare, stimolazione culturale e supporto
    educativo possono contribuire allo sviluppo delle capacità cognitive.
\end{snippet}

\begin{snippet}{intelligenza-eta}
    Negli anziani in buona salute vi sono poche prove a favore di un declino
    cognitivo generalizzato. L'intelligenza cristallizzata tende a conservarsi
    meglio, mentre possono diminuire alcune funzioni come memoria di lavoro,
    memoria a lungo termine e velocità di elaborazione.
\end{snippet}

\begin{snippetdefinition}{riserva-cognitiva-definition}{Riserva cognitiva}
    La \emph{riserva cognitiva} è un insieme di risorse a livello neurale che
    permette di far fronte in modo più efficace a eventi traumatici o al declino
    cognitivo.
\end{snippetdefinition}

\begin{snippet}{riserva-cognitiva-fattori}
    Attività sociali, fisiche e intellettuali svolte nel corso della vita sono
    associate a una migliore conservazione delle capacità cognitive. Questionari
    come il \emph{CRIq} stimano la riserva cognitiva considerando lavoro
    pregresso, attività del tempo libero e attività ancora praticate.
\end{snippet}

\begin{snippetdefinition}{indice-ereditabilita-definition}{Indice di ereditabilità}
    L'\emph{indice di ereditabilità} fornisce una stima del grado in cui la
    variabilità di una caratteristica in una popolazione è determinata dalla
    variabilità genetica.
\end{snippetdefinition}

\begin{snippet}{genetica-ambiente-intelligenza}
    Gli studi sui gemelli mostrano che il \(QI\) è influenzato in modo
    importante dalla componente genetica. Tuttavia, l'ambiente condiviso e le
    condizioni educative modulano tale influenza: la somiglianza genetica non
    basta da sola a spiegare tutte le differenze nei punteggi di intelligenza.
\end{snippet}

\includesnpt[width=62\%,src=/snippet/static-psychology/iq-genetics-environment-correlations.jpg]{centered-img}

\begin{snippet}{intelligenza-conclusione}
    I dati disponibili indicano che l'intelligenza dipende dall'interazione tra
    componenti genetiche, sviluppo biologico, esperienza, ambiente educativo e
    opportunità sociali. La componente genetica è rilevante, ma l'ambiente ha un
    ruolo non trascurabile.
\end{snippet}

\end{document}
